\subsection{Diversity at finite steps and sampling budgets}
\label{app:finite-time-diversity}

Softmax outcome logits amplify probability differences among equally rewarded
solutions \citep[Theorem~3.1 and App.~A.1]{sinha2026expected}, and
sampling imbalances can reinforce this effect
\citep[Section~4.2]{lochab2026ucpo}. Under the categorical model,
\pmd{} decreases from every nonuniform correct-mode distribution along
both the mean flow and finite expected-gradient steps. Sampled updates can
also break an initially uniform distribution's symmetry. The bounds below
relate concentration to the probability of observing a competing solution
mode within a finite inference budget.

\begin{proposition}[Exact diversity drift and matched correctness]
\enforcehalffinalproseline
\label{prop:pcmd-drift}
Under Assumption~\ref{ass:categorical-mean-flow}, consider
$\dot z=c(P)\nabla_zP$ with $c(P)>0$. Write
$D(q)=1-\sum_cq_c^2$, $S_2=\sum_cq_c^2$, and
$V(q)=\sum_cq_c^3-S_2^2=\operatorname{Var}_{C\sim q}(q_C)$.
For incorrect conditional probabilities $r_a=p_a/(1-P)$, set
$T_2=\sum_ar_a^2$. In effective time $d\tau=c(P)P(1-P)dt$,
\begin{equation}
 \frac{dD}{d\tau}=-2V(q),
 \qquad
 \frac{dD}{d\log[P/(1-P)]}
 =-\frac{2V(q)}{S_2+T_2}.
 \label{eq:pcmd-drift}
\end{equation}
Thus diversity strictly decreases at every nonuniform interior $q$.
Positive-coefficient mean flows from the same initial logits follow the same
$D$-versus-$P$ curve, independently of their chosen outcome-binary advantage or
group size.
\end{proposition}

\begin{proof}
\enforcehalffinalproseline
For the replicator flow~\eqref{eq:grpo-replicator}, the potential
$U(q)=S_2/2$ has fitness $\partial U/\partial q_c=q_c$.
The fitness-variance identity of
\citet[Theorems~1--2]{harper2009informationgeometry} gives
$dU/d\tau=V(q)$, hence $dD/d\tau=-2V(q)$.
The variance vanishes exactly when the
positive coordinates of $q$ are equal. Also,
\[
 \dot P=c(P)\|\nabla_zP\|_2^2
       =c(P)P^2(1-P)^2(S_2+T_2),
\]
so $d\log[P/(1-P)]/d\tau=S_2+T_2>0$, proving the second identity.
Finally, the full logit equation in effective time is
$dz_c/d\tau=q_c$ and $dz_a/d\tau=-r_a$; it contains no $c$.
The same initial logits therefore give the same full trajectory after the
time change, and hence the same curve at matched correctness for every positive-coefficient outcome-binary estimator considered here.
\end{proof}

\begin{theorem}[Collapse under finite expected-gradient steps]
\enforcehalffinalproseline
\label{thm:finite-expected-steps}
Retain (A1), (A2), (A4), and (A5) of
Assumption~\ref{ass:categorical-mean-flow}, and replace (A3) by
\begin{equation}
 z_{k+1}=z_k+\eta_k c(P_k)\nabla_zP_k,
 \qquad 0<\eta_k<\infty,\quad \sum_{k=0}^\infty\eta_k=\infty,
 \label{eq:finite-expected-update}
\end{equation}
where $c$ is continuous and positive on $(0,1)$. Each step strictly decreases
$D(q_k)$ unless $q_k$ is uniform. If $q_0$ has a unique largest coordinate
$c^\star$, then $P_k\to1$ and $q_{c^\star,k}\to1$.
No upper bound on the finite step sizes is required.
\end{theorem}

\begin{proof}
\enforcehalffinalproseline
Let $h_k=\eta_kc(P_k)P_k(1-P_k)>0$. The exact updates are
$z_c^+=z_c+h_kq_c$ and $z_a^+=z_a-h_kr_a$, so
\begin{equation}
 q_c^+=\frac{q_c e^{h_kq_c}}{\sum_dq_d e^{h_kq_d}}.
 \label{eq:finite-correct-tilt}
\end{equation}
For fixed $q$, put $u_c(s)=q_ce^{sq_c}/\sum_dq_de^{sq_d}$.
Differentiation gives
\[
 \frac{d}{ds}\sum_cu_c(s)^2
 =\sum_{c,d}u_c(s)u_d(s)
       \big(u_c(s)-u_d(s)\big)(q_c-q_d)\ge0.
\]
The ordering of $u(s)$ is that of $q$, and the inequality is strict whenever
$q$ is nonuniform. Integrating from $0$ to $h_k$ proves the diversity claim.
Equation~\eqref{eq:finite-correct-tilt} also preserves the leader and makes
its conditional probability nondecreasing.

Every correct logit increases and every incorrect logit decreases, hence
$P_k$ increases. Jensen's inequality gives
\[
 \log\frac{P_{k+1}}{1-P_{k+1}}-\log\frac{P_k}{1-P_k}
 =\log\sum_cq_c e^{h_kq_c}-\log\sum_ar_a e^{-h_kr_a}
 \ge h_k\sum_cq_c^2\ge h_k/m.
\]
If $P_k$ had an interior limit, continuity and positivity of $c$ would imply
$\sum_kh_k=\infty$, contradicting bounded correctness odds. Thus $P_k\to1$.
Moreover, $\sum_kh_k=\infty$ is necessary: otherwise every logit would have
finite total variation and converge to a finite value, contradicting
$P_k\to1$.
Set $\alpha=q_{c^\star,0}$ and
$\xi_{c,k}=q_{c,k}/q_{c^\star,k}$. For $c\ne c^\star$,
\[
 \xi_{c,k+1}=\xi_{c,k}
  e^{-h_k(q_{c^\star,k}-q_{c,k})},
 \qquad
 q_{c^\star,k}-q_{c,k}
 \ge\alpha(1-\xi_{c,0})>0.
\]
The divergent sum of $h_k$ forces every such ratio to zero, proving
$q_{c^\star,k}\to1$.
\end{proof}

The update uses the exact expected gradient. Finite steps preserve the
concentration conclusion but can change the trajectory across estimators.
The identical diversity-versus-correctness curve in
Proposition~\ref{prop:pcmd-drift} requires infinitesimal mean flow;
it does not extend to these finite updates.

Finite multinomial counts can break a uniform policy's symmetry
\citep[Proposition~A.1]{lochab2026ucpo}. For the centered group estimator,
the leading expected one-step loss in \pmd{} is quadratic in the learning
rate, with a coefficient determined by group size and correctness.

\begin{proposition}[Sampled updates break conditional symmetry]
\enforcehalffinalproseline
\label{prop:stochastic-pcmd}
Retain (A1), (A2), (A4), and (A5) of
Assumption~\ref{ass:categorical-mean-flow}, and replace (A3) by one sampled
update $z^+=z+\eta\widehat g_G$ from
Equation~\eqref{eq:group-estimator}, with $\eta>0$. Suppose $0<P<1$ and $q_c=1/m$ for every correct mode. Then
\begin{equation}
 \E[D(q^+)]<1-1/m,
 \qquad \E[q_c^+]=1/m.
 \label{eq:stochastic-strict-decrease}
\end{equation}
As $\eta\to0$, for the fixed initial policy and estimator,
\begin{equation}
 \E[D(q^+)]
 =1-\frac1m
  -\eta^2\frac{m-1}{m^3}
    \E\!\left[\frac{w(R)^2R(G-R)^2}{G^4}\right]
  +O(\eta^3),
 \label{eq:stochastic-pcmd-expansion}
\end{equation}
where the summand is zero on $R=0,G$. For Dr.GRPO, $w=1$ and
\begin{equation}
 \E\!\left[\frac{R(G-R)^2}{G^4}\right]
 =\frac{G-1}{G^3}P(1-P)\big[1+(G-2)(1-P)\big].
 \label{eq:stochastic-dr-coefficient}
\end{equation}
\end{proposition}

\begin{proof}
\enforcehalffinalproseline
Let $N_c$ count occurrences of correct mode $c$. The centered advantages
sum to zero, so the common softmax-score term cancels and
$\widehat g_c=\kappa_RN_c$, where
$\kappa_R=w(R)(G-R)/G^2$ on mixed groups and $\kappa_R=0$ otherwise.
Conditional on $R$, the correct counts are
$N\sim\operatorname{Multinomial}(R;1/m,\ldots,1/m)$, and
\[
 q_c^+=\frac{e^{\eta\kappa_RN_c}}{\sum_de^{\eta\kappa_RN_d}}.
\]
Diversity is at most $1-1/m$, with equality only at uniform $q^+$.
The event $R=1$ has positive probability; on this event one correct logit
increases and the others do not, so $q^+$ is nonuniform. This proves the
strict inequality for every $\eta>0$. Permutation symmetry of the correct
labels gives $\E[q_c^+]=1/m$.

For a fixed vector $v$, Taylor expansion about uniform logits gives
\[
 \sum_c\operatorname{softmax}(\eta v)_c^2
 =\frac1m+\frac{\eta^2}{m^2}\sum_c(v_c-\bar v)^2+O(\eta^3),
 \qquad \bar v=\frac1m\sum_cv_c.
\]
Use $v_c=\kappa_RN_c$ and
$\E[\sum_c(N_c-R/m)^2\mid R]=R(1-1/m)$.
There are finitely many possible groups and all mixed-group weights are
finite, so the Taylor remainder is uniform over those groups. Averaging
proves Equation~\eqref{eq:stochastic-pcmd-expansion}.
Finally,
\[
 \E[R(G-R)^2]
 =G(G-1)(G-2)P(1-P)^2+G(G-1)P(1-P),
\]
by the binomial factorial moments, yielding
Equation~\eqref{eq:stochastic-dr-coefficient}.
\end{proof}

At fixed $P\in(0,1)$, the Dr.GRPO coefficient in
Equation~\eqref{eq:stochastic-dr-coefficient} is
$P(1-P)^2/G+O(G^{-2})$ as $G\to\infty$. The coefficient of the
quadratic learning-rate term therefore decays asymptotically as $1/G$;
it need not decrease at every increment in finite group size. This sampling
effect is absent from the uniform mean flow. The result concerns one step
from a uniform correct-mode distribution and gives neither monotone expected
diversity from an arbitrary initial policy nor stochastic-optimizer convergence.

\begin{corollary}[Finite-budget disappearance of competing modes]
\enforcehalffinalproseline
\label{cor:finite-budget-invisibility}
In the mean flow of Theorem~\ref{thm:grpo-collapse}, let $c^\star$ be the
unique initial correct leader and define
\[
 \alpha=q_{c^\star}(0),
 \qquad
 \Gamma_0=\sum_{c\ne c^\star}
             \frac{q_c(0)}{\alpha-q_c(0)}.
\]
At effective time $\tau$,
\begin{equation}
 1-q_{c^\star}(\tau)\le\Gamma_0e^{-\alpha\tau},
 \qquad D(q(\tau))\le2\Gamma_0e^{-\alpha\tau}.
 \label{eq:finite-minority-bound}
\end{equation}
For $K\ge1$ independent full-policy draws at that checkpoint,
\begin{equation}
 \Pr(\text{any correct mode other than }c^\star\text{ is observed})
 \le K\Gamma_0e^{-\alpha\tau}.
 \label{eq:finite-inference-visibility}
\end{equation}
In particular, if
$\tau\ge\max\{0,\alpha^{-1}\log(K\Gamma_0/\delta)\}$ for
$0<\delta<1$, no competing correct mode is observed with probability at
least $1-\delta$, although every mode retains positive probability.
\end{corollary}

\begin{proof}
\enforcehalffinalproseline
For $c\ne c^\star$, write $\xi_c=q_c/q_{c^\star}<1$.
The leader is nondecreasing, and Equation~\eqref{eq:grpo-replicator} gives
\[
 \frac{d}{d\tau}\log\frac{\xi_c}{1-\xi_c}
 =-q_{c^\star}\le-\alpha.
\]
Hence
$\xi_c(\tau)\le[q_c(0)/(\alpha-q_c(0))]e^{-\alpha\tau}$.
Now $1-q_{c^\star}=\sum_{c\ne c^\star}\xi_c/(1+\sum_{c\ne c^\star}\xi_c)$
and $D\le1-q_{c^\star}^2\le2(1-q_{c^\star})$, proving the first bounds.
A single full-policy draw hits a competing correct mode with probability
$P(1-q_{c^\star})\le1-q_{c^\star}$; a union bound over $K$ draws completes
the proof.
\end{proof}

The bound gives a finite-budget interpretation of the rare-output pruning
mechanism discussed by \citet[Section~4.2.5]{lochab2026ucpo}.
It can also be indexed by correctness: Proposition~\ref{prop:pcmd-drift}
and $S_2+T_2\le2$ imply
$\tau\ge\tfrac12\log\{[P/(1-P)]/[P(0)/(1-P(0))]\}$.
These are sufficient visibility bounds within the categorical model;
they do not require literal support loss or identify an empirical collapse
time for neural training.


